% ROUND63 OE manuscript — working skeleton (pre-F1 draft)
% PORTING NOTE: written against the standard article class so it compiles
% anywhere; port to opticajnl.cls (Optica template) at submission prep — the
% section structure below already follows the OE research-article shape.
% PLACEHOLDER DISCIPLINE: every number that depends on S1/S2/S3/S4 outputs is
% written as \SPH{...} (S-placeholder) and indexed in PLACEHOLDER_LEDGER.md;
% the manuscript is "filled" only after the confirmatory runs, and zero \SPH
% may survive submission (campaign rule inherited from the 2026-07-13 audit).
\documentclass[10pt,twocolumn]{article}
\usepackage{amsmath,amssymb,bm}
\usepackage{graphicx}
\usepackage[margin=2.2cm]{geometry}
\usepackage{xcolor}
\usepackage{hyperref}
\newcommand{\SPH}[1]{\textcolor{red}{[#1]}}   % unfilled campaign placeholder
\newcommand{\rhobar}{\bar{\rho}}

\title{Dead-time-aware high-flux single-pixel imaging: operating beyond the
conventional photon-counting regime with renewal-statistical reconstruction}

% Author block: USER DECISION (affiliation + collaborators) — do not fill.
\author{\SPH{author block — user decision}}
\date{}

\begin{document}
\maketitle

\begin{abstract}
% Chen-genre shape: field-first opening, pain point, "Here we report ...",
% mechanism, headline numbers, application close. NO "for the first time"
% (round-4 novelty discipline), no "first".
Single-pixel imaging with photon-counting detectors is conventionally operated
at low photon flux, keeping the detector dead-time loss below a few percent at
the cost of long acquisition times. Here we report dead-time-aware high-flux
single-pixel imaging, in which the operating flux is deliberately raised into
the strongly nonlinear regime $\rhobar \in [0.3, 1]$ (dead-time duty $\rhobar
= \lambda\tau$) and the nonlinearity is absorbed by a convex renewal
quasi-likelihood (RQL) reconstruction with a deployment-legal, truth-free
regularization rule. A finite-window count-information analysis explains why
this is possible and where it must stop: the Fisher information of the
integrated photon count possesses a principal ridge at the cube-root load
$\rho^{*}(\nu) = (6\nu)^{1/3} - 2/3 + O(\nu^{-1/3})$ ($\nu = T/\tau$ dead-time
slots per exposure), arising from terminal detector-phase information loss.
Preregistered simulations across \SPH{S2-A grid summary} show a median
time-to-quality acceleration of \SPH{S2 median $S$ with CI} at matched
radiometric image quality relative to the conventional $\rhobar = 0.05$
operating point, robust to \SPH{S3 mismatch envelope summary}. The proposed
operating-point map opens up an avenue for photon-budget-limited and
dwell-time-limited single-pixel systems.
\end{abstract}

\section{Introduction}
% Para 1 — field-first (Chen style): SPI/ghost imaging + photon counting
% context, applications (low-light, NIR/SWIR, LiDAR-adjacent, biological).
\SPH{field-first paragraph}

% Para 2 — the pain point: the 5\% rule / low-flux convention; dead time as
% the reason; time cost scales inversely with allowed flux; cite Liu-2021-type
% pre-correction, Li-2023 RJMCMC, Jorgensen-2026 dead-time MLE as the
% correction-centric line (defensive citations, round-4 list).
\SPH{pain-point paragraph}

% Para 3 — the reframe (our contribution): not "correct the distortion" but
% "move the operating point"; RQL joint inverse problem; the ridge as the
% principled boundary. Safe novelty wording VERBATIM from round-4/5 ruling:
Previous studies established count-rate-dependent CRLB degradation,
task-specific optimal operating fluxes, and high-flux recovery using
time-resolved dead-time models \SPH{cites: Wang 2011; B\'ecares--Bl\'azquez
2012; Gupta et al.; Rapp \& Goyal; Jorgensen--Johnson 2026; 2025 dToF}. Here,
we study a different information bottleneck: a finite exposure from which only
the integrated nonparalyzable count is retained. For this count-only
observation, we derive a principal finite-window Fisher-information ridge and
show that its optimal load obeys a cube-root law, $\rho^{*}(\nu) =
(6\nu)^{1/3} - 2/3 + O(\nu^{-1/3})$, arising from terminal detector-phase
information loss.

% Para 4 — contributions list + preregistration statement + paper roadmap.
\SPH{contributions + prereg-statement paragraph}

\section{Physical model and problem statement}
\label{sec:model}
% Frozen content (spec D2.1 §2): non-paralyzable renewal counting, active
% start, exact PMF; dimensionless axes; arrival rates lambda_i = Phi <a_i, x>
% + d; raw-count observation; the exact PMF via the Poisson representation.
A single-pixel detector with dead time $\tau$ observes, for illumination
pattern $a_i$ and scene $x \ge 0$, a Poisson arrival process of rate
$\lambda_i = \Phi\, a_i^{\top} x + d$ during an exposure of duration
$T = \nu\tau$; after each registered count the detector is blind for $\tau$
(non-paralyzable, active start). The recorded observation is the integrated
count $N_i$ alone, with exact law
\begin{equation}
  P(N_i \ge m) = F_{\Gamma(m,\lambda_i)}\!\bigl(T - (m-1)\tau\bigr)
              = P\bigl(\mathrm{Pois}(z_m) \ge m\bigr),
  \quad z_m = \lambda_i\,(T-(m-1)\tau).
  \label{eq:exactpmf}
\end{equation}
The per-slot load $\rho_i = \lambda_i \tau$ is the dead-time duty; the
conventional regime keeps $\rho \lesssim 0.05$.
\SPH{paragraph: measurement ensembles (Bernoulli, complementary Hadamard
pairs at $2M$ exposures, gamma-modulated), per-pattern $\rho$ percentiles}

\section{Renewal quasi-likelihood reconstruction}
\label{sec:rql}
% Frozen: the convex objective; why not full Gaussian (supplement pointer);
% relation of the single-frame root to pre-correction (honesty clause).
We reconstruct by
\begin{equation}
  \hat{x} = \arg\min_{x \ge 0}\;
  \frac{1}{M}\sum_{i=1}^{M}\Bigl[(T - N_i\tau)\,\lambda_i(x)
  - N_i \log \lambda_i(x)\Bigr] + \lambda_{\mathrm{TV}}\,
  \mathrm{TV}(x),
  \label{eq:rql}
\end{equation}
with $\lambda_i(x) = \Phi a_i^{\top} x + d$. The data term is convex in
$\lambda$ (and in $x$ through the affine map); its single-measurement
stationary point $\hat\lambda = N/(T - N\tau)$ coincides with the classical
dead-time rate correction, so the estimator's novelty resides not in the
per-measurement formula but in the joint regularized inverse problem and the
operating-point map it enables. At ceiling counts ($N\tau \ge T$) the
objective is unbounded below in $\lambda$ — the true identifiability boundary
of count-only observation; the fraction of ceiling counts is recorded per
cell. A full heteroscedastic Gaussian likelihood fails at high load through a
variance-collapse non-coercivity (Supplement~\SPH{S-ref}).

\subsection{Truth-free regularization selection and descriptive audit}
\label{sec:selection}
% Frozen F1 rule (round-4 + round-5): DEV/AUDIT split, DEV-only one-SE
% selection, descriptive audit. Prereg wording VERBATIM (round-5):
For each measurement set, patterns are split once into a development part
(80\% of pattern groups; complementary pairs kept whole) and an audit part
(20\%). The TV weight is chosen on the development part alone: a grouped
5-fold cross-fit scores each candidate $\eta = \lambda_{\mathrm{TV}} /
\lambda_{\mathrm{max}}$ on a dimensionless 9-point path by the calibrated
renewal residual $r_i = (N_i - \mu_{\mathrm{NP}}(\hat\lambda_i)) /
\sqrt{V_{\mathrm{NP}}(\hat\lambda_i)}$, and the largest $\eta$ within one
fold-dispersion standard error of the minimizer is frozen. All arms are
selected by the identical rule.

The measurement audit is descriptive and does not constitute a model-adequacy
test. All audit statistics are excluded from regularization selection,
reconstruction, cell inclusion, and confirmatory pass/fail decisions. The
refit-bootstrap reference distribution is conditional on a smoothed plug-in
scene; its empirical tail ranks are therefore not calibrated model-class
$p$-values for arbitrary scenes and can be anti-conservative in the upper
tail when the true scene is rougher than the plug-in. Moreover, count-only
observations collected at a single operating point cannot generally
distinguish detector mean-response mismatch from radiometric rescaling of the
scene. Detector-mismatch consequences are therefore quantified directly
through the preregistered radiometric metrics in
Section~\ref{sec:mismatch}.

\section{The finite-window count-information ridge}
\label{sec:ridge}
% Port from paper/THEORY_SECTION_DRAFT.md (X.1-X.4) — the crown jewel.
\SPH{port theory section: missing-information identity, cube-root law,
peak-information expansion, VarN-still-diverges remark, zoning + CLT
boundary; Fig.~\ref{fig:main}(a)}

\section{Preregistered simulation campaign}
\label{sec:campaign}
% Frozen: prereg hypothesis verbatim (D2.1 §0), grids, metrics, S_j
% definition, bootstrap, freeze/commit provenance paragraph.
We test whether renewal quasi-likelihood reconstruction permits operation in
$\rhobar \in [0.3, 1]$ with a statistically significant reduction in total
optical integration time relative to the conventional $\rhobar = 0.05$
operating point. \SPH{campaign-design paragraph: grid, images-as-units,
time-to-quality protocol (fixed $M$, dwell scan), radiometric PSNR primary,
censoring, stratified bootstrap, S1/S2 separation, freeze commit SHA}

\section{Results}
\label{sec:results}
\subsection{Time-to-quality acceleration at high flux}
\SPH{S2-A headline: median $S_j$, bootstrap CI, per-image spread;
Fig.~\ref{fig:main}(b,c)}
\subsection{Operating-point map and the ridge}
\SPH{empirical $S_Q$ vs $\rhobar$ against the exact-FI envelope;
Fig.~\ref{fig:main}(d); $\rhobar=2$ information-decreasing confirmation}
\subsection{Robustness to detector mismatch}
\label{sec:mismatch}
\SPH{S3: tau error, dark, afterpulsing, start modes, guard; radiometric
degradation envelopes; descriptive audit distributions (supplement)}
\subsection{Pattern ensembles, undersampling, and scale}
\SPH{S2-B/C: M/n, hadpair (2M exposure accounting), gam4, 128\textsuperscript{2}}
\subsection{Exact-likelihood reference at small scale}
\SPH{S4: scalar Fisher map + 8/16 px exact-vs-RQL split}

\begin{figure*}[t]
  \centering
  \SPH{4-panel main figure per spec D2.1 §6 — (a) exists:
  results/round63\_theory/fig\_a\_ridge\_map.pdf}
  \caption{\SPH{main figure caption}}
  \label{fig:main}
\end{figure*}

\section{Discussion}
% Honest scoping: simulation-only (compensated by detector-parameter ranges +
% full mismatch map + open frozen benchmark); narrow novelty scope of the
% ridge (round-4 wording); count-only vs timestamp architectures (Rapp);
% "opens up an avenue" close (Chen genre). NO separate limitations section
% per Chen style, but the honesty lives inside these paragraphs.
\SPH{discussion paragraphs}

\section*{Code and data availability}
All simulation code, frozen preregistration documents, per-cell manifests,
and evidence bundles are available at \SPH{repo URL wording — user decision}.

\section*{Acknowledgments}
\SPH{funding/acknowledgments — user decision}

% \bibliography{refs}
\SPH{bibliography — build refs.bib during S5 with defensive citations from
docs/ROUND63\_GPT\_ROUND4\_DIGEST.md \S Q2}

\end{document}
